\documentclass[11pt]{article}
\usepackage{ap-physc}
\course{AP Physics C: Electricity and Magnetism}
\headcourse{AP Physics C: E\&M}
\headpart{Unit 10 \textperiodcentered\ Question Set}
\title{Unit 10: Conductors and Capacitors: Question Set}
\author{Paul Wang}
\date{November 2026}
\docpart{Unit 10 \textperiodcentered\ Question Set}

\begin{document}
\maketitle

\begin{abstract}
Exam-style practice on Unit 10 (topics 10.1--10.4), which is 10--15\% of the multiple-choice section. Section~1--3 hold 12 four-choice multiple-choice questions grouped by topic; most need a short calculation or a line of reasoning, as on the exam. The last page before the key is a full Translation Between Representations free-response question (12 points). All items are original and written for practice. Use the AP equation sheet and a calculator. Allow about 24 minutes for the multiple choice (the exam's pace is about 2 minutes per question) and 24 minutes for the free response. The answer key explains every item and names its topic; a scoring guide for the free response follows it.
\end{abstract}

\scoringline{MCQ: \blank\ / \qtotal \qquad FRQ: \blank\ / 12 \qquad Target: 8+ MCQ \qquad Pace: 2 min per MCQ}

\section{Conductors (10.1--10.2)}

\begin{q}{2}{10.1 Conducting shell}{A Gaussian sphere inside the metal has $E = 0$, so it encloses zero net charge: the inner surface carries $-q$. The shell is neutral, so $+q$ goes to the outer surface.}
A point charge $+q$ sits at the center of a thick, neutral, conducting spherical shell. What charges appear on the shell's inner and outer surfaces?
\chh{Inner: $0$; outer: $0$}{Inner: $-q$; outer: $+q$}{Inner: $+q$; outer: $-q$}{Inner: $-q$; outer: $0$}
\end{q}
\begin{q}{4}{10.1 Field at a conductor's surface}{A pillbox with one face inside the metal ($E = 0$) and one just outside gives $EA = \sigma A/\varepsilon_0$. Any parallel component would move surface charges, so the field is perpendicular.}
A conductor's surface carries charge density $\sigma$ at some point. What is the field just outside the surface there?
\ch{$\sigma/(2\varepsilon_0)$, perpendicular to the surface}{$\sigma/\varepsilon_0$, parallel to the surface}{$0$}{$\sigma/\varepsilon_0$, perpendicular to the surface}
\end{q}
\begin{q}{1}{10.1 Potential inside a conductor}{$E = 0$ inside, so $V$ does not change from the surface inward: it stays $kQ/R$.}
A metal sphere of radius $R$ carries charge $Q$. What is the electric potential at a point a distance $r < R$ from its center?
\chh{$kQ/R$}{$0$}{$kQ/r$}{$kQr/R^3$}
\end{q}
\begin{q}{3}{10.2 Sharing charge}{Equal potentials make $q \propto R$: $q_1/q_2 = 1/3$, so $q_1 = \tfrac14(8.0) = 2.0\ \mu\mathrm{C}$.}
Metal spheres of radii 0.10 m and 0.30 m, far apart, are joined by a thin wire and share a total charge of $8.0\ \mu\mathrm{C}$. What charge ends up on the smaller sphere?
\chh{$6.0\ \mu\mathrm{C}$}{$4.0\ \mu\mathrm{C}$}{$2.0\ \mu\mathrm{C}$}{$0.80\ \mu\mathrm{C}$}
\end{q}
\begin{q}{3}{10.2 Surface field}{$E_{surface} = kq/R^2$ and $q \propto R$, so $E \propto 1/R$: three times stronger on the small sphere.}
For the two joined spheres in the previous question, which has the stronger electric field just outside its surface?
\ch{The larger sphere}{They are equal, because they are at the same potential}{The smaller sphere}{Neither; the field outside a conductor is zero}
\end{q}

\section{Capacitors (10.3)}

\begin{q}{2}{10.3 Parallel-plate capacitance}{$C = \varepsilon_0A/d = (8.85\times10^{-12})(0.020)/(1.0\times10^{-3}) = 1.8\times10^{-10}$ F $= 180$ pF.}
A parallel-plate capacitor has plates of area $0.020\ \mathrm{m^2}$ separated by 1.0 mm of air. What is its capacitance? ($\varepsilon_0 = 8.85\times10^{-12}\ \mathrm{C^2/(N\,m^2)}$)
\chh{$1800$ pF}{$180$ pF}{$18$ pF}{$1.8$ pF}
\end{q}
\begin{q}{4}{10.3 Spherical capacitor}{With $\pm Q$, $E = kQ/r^2$ between the spheres, so $\Delta V = kQ\left(\dfrac1a - \dfrac1b\right) = kQ\dfrac{b - a}{ab}$ and $C = Q/\Delta V$.}
A spherical capacitor has an inner conductor of radius $a$ and an outer shell of radius $b$. What is its capacitance?
\ch{$4\pi\varepsilon_0(b - a)$}{$\dfrac{\varepsilon_0 4\pi b^2}{b - a}$}{$4\pi\varepsilon_0\dfrac{b - a}{ab}$}{$4\pi\varepsilon_0\dfrac{ab}{b - a}$}
\end{q}
\begin{q}{2}{10.3 Combining capacitors}{Parallel pair: $6.0\ \mu\mathrm{F}$. In series with 3.0: $\dfrac{(6.0)(3.0)}{9.0} = 2.0\ \mu\mathrm{F}$.}
A $4.0\ \mu\mathrm{F}$ and a $2.0\ \mu\mathrm{F}$ capacitor are in parallel, and that pair is in series with a $3.0\ \mu\mathrm{F}$ capacitor. What is the equivalent capacitance?
\chh{$1.3\ \mu\mathrm{F}$}{$2.0\ \mu\mathrm{F}$}{$4.5\ \mu\mathrm{F}$}{$9.0\ \mu\mathrm{F}$}
\end{q}
\begin{q}{1}{10.3 Energy when the gap changes}{Isolated, so $Q$ is fixed; $C$ halves; $U = Q^2/(2C)$ doubles. The extra energy is the work done pulling the attracting plates apart.}
A charged parallel-plate capacitor is disconnected from its battery, and then its plate separation is doubled. What happens to its stored energy?
\chh{It doubles.}{It halves.}{It is unchanged.}{It quadruples.}
\end{q}
\begin{q}{3}{10.3 Energy density}{$u = \tfrac12\varepsilon_0E^2 = \tfrac12(8.85\times10^{-12})(10^{12}) = 4.4\ \mathrm{J/m^3}$.}
The field between a capacitor's plates is $1.0\times10^6$ V/m (air). What is the energy density of the field?
\chh{$4.4\times10^{-6}\ \mathrm{J/m^3}$}{$0.50\ \mathrm{J/m^3}$}{$4.4\ \mathrm{J/m^3}$}{$8.9\ \mathrm{J/m^3}$}
\end{q}

\section{Dielectrics (10.4)}

\begin{q}{4}{10.4 Dielectric with battery connected}{The battery holds $V$ fixed. $C$ triples, so $Q = CV$ triples; the battery supplies the extra charge.}
A capacitor stays connected to a battery while a slab of dielectric constant 3 is slid in to fill the gap. Which is true afterward?
\ch{The voltage drops to one third and the charge is unchanged.}{Both the charge and the voltage triple.}{The charge drops to one third and the voltage is unchanged.}{The charge triples and the voltage is unchanged.}
\end{q}
\begin{q}{1}{10.4 Partially filled capacitor}{Each half has area $A/2$ and the full gap, so the halves are in parallel: $C = \tfrac12 C_0 + \tfrac12\kappa C_0 = 2C_0$.}
A slab of dielectric constant $\kappa = 3$ fills the left half of the area between a parallel-plate capacitor's plates, across the full gap. The empty capacitor had capacitance $C_0$. What is the new capacitance?
\chh{$2C_0$}{$3C_0$}{$1.5C_0$}{$4C_0$}
\end{q}

\section{Free response}

\begin{frq}{Question 1 \textperiodcentered\ Translation Between Representations \textperiodcentered\ 12 points \textperiodcentered\ 24 minutes}
An air-filled parallel-plate capacitor with plate area $A$ and gap $d$ is charged by a battery of voltage $V_0$. The battery is then disconnected. An insulating handle is used to pull the plates apart slowly, from separation $d$ to $2d$. Ignore edge effects.

\begin{enumerate}[label=(\alph*),leftmargin=2em,itemsep=4pt]
\item Draw the capacitor at separation $d$, showing the sign of the charge on each plate and the electric field lines between and outside the plates. \pts{2}
\item On three separate sets of axes, sketch graphs of the charge $Q$, the voltage $\Delta V$, and the field magnitude $E$ between the plates, each as a function of the separation $x$ for $d \le x \le 2d$. Label the starting and ending values in terms of $V_0$ and the starting values. \pts{4}
\item Derive an expression for the work done by the person pulling the plates from $d$ to $2d$, in terms of $A$, $d$, $V_0$, and physical constants. \pts{3}
\item Suppose instead the battery stays connected while the plates are pulled from $d$ to $2d$. Describe how each of your three graphs in (b) would change, and state whether the capacitor's stored energy increases or decreases. Justify. \pts{3}
\end{enumerate}
\end{frq}

\makekey

\subsection*{Free-response scoring guide (12 points)}
\small
Award each point only when the answer does what the model does. Points are independent unless a part says otherwise; a correct method with an answer carried forward from an earlier error earns the later points.
\begin{scoring}
(a) (2) & 1 pt: $+$ on one plate, $-$ on the other. 1 pt: evenly spaced parallel lines from $+$ to $-$ between the plates, and no field outside. & Field lines that curve or bunch between the plates, or that extend outside. \\
(b) (4) & 1 pt: $Q$ is a horizontal line at $Q_0 = \varepsilon_0AV_0/d$ (isolated). 1 pt: $\Delta V$ rises linearly from $V_0$ to $2V_0$ ($\Delta V = Q_0x/(\varepsilon_0A)$). 1 pt: $E$ is a horizontal line at $V_0/d$ ($E = \sigma/\varepsilon_0$ is set by the fixed charge). 1 pt: all endpoints labeled correctly. & Drawing $\Delta V$ constant, as if the battery were still connected. \\
(c) (3) & 1 pt: $W = \Delta U$ (slow pull, no kinetic energy). 1 pt: $U_1 = \tfrac12 C_0V_0^2$, $U_2 = \dfrac{Q_0^2}{2(C_0/2)} = C_0V_0^2$. 1 pt: $W = \tfrac12 C_0V_0^2 = \dfrac{\varepsilon_0AV_0^2}{2d}$. (Equivalently: constant force $Q_0E/2$ on a plate times distance $d$.) & Using $U = \tfrac12 CV^2$ with $V$ held at $V_0$, which gives a decrease in energy. \\
(d) (3) & 1 pt: $\Delta V$ is now constant at $V_0$. 1 pt: $Q = CV_0$ falls to $Q_0/2$ (as $1/x$), and $E = V_0/x$ falls to $V_0/(2d)$. 1 pt: stored energy $\tfrac12 CV_0^2$ decreases (to half), because $C$ halves at fixed voltage; charge flows back into the battery. & Saying the energy increases because work is done on the plates; the battery absorbs energy as well. \\
\end{scoring}

\end{document}
